\section{Controlled Replay Details and Revised Accounting}\label{app:controlled-details}
The current experimental source is \path{experiments/reviewer-audit/recompute.py}, with the instrumented retail implementation in \path{retail_replay.py}. Historical inputs and previous summaries are retained unchanged; revised JSON files take precedence for the current manuscript. No fresh model generations were requested.

\paragraph{TLS corpus and length accounting}
The 48 saved outputs comprise 24 brief/expanded pairs. Training and testing each contain 12 topics and 24 outputs. T1 uses the original verbosity labels. T2 retains brief responses and takes the first two period-delimited sentences of expanded responses. T3 uses the first two and first three sentences of each expanded response. T4 retains both original outputs per topic and labels their domain. There is one generation per prompt, not repeated model sampling. T4 has 14 positive and 10 negative test outputs (seven versus five topics), so 50\% is not its empirical majority-class baseline.

For every regime and framing policy, \path{tls_records.json} saves topic, split, label, plaintext bytes, application-record count, payload bytes, wire bytes and recovered plaintext hash. \path{tls_gradient.json} reports absolute means and matched denominators. All 48 texts are used for cost means; only 24 held-out outputs are used for accuracy. A record contains five TLS header bytes plus 24 encryption-overhead bytes under the pinned TLS~1.2 suite. Handshake traffic, TCP/IP and timing are outside this cost scope. Bucket 1024 uses four 256-byte writes per quantum, hence 1,140 rather than 1,082 wire bytes per quantum. Buckets 256/512 use 128-byte writes.

The training-fitted cap rounds the maximum framed training length, 1,194 bytes, up to 1,280 bytes. It emits ten records of 128 bytes (1,570 wire bytes). Overflow adds 128-byte frames, never truncates, and leaks an overflow-dependent length; none of the saved test responses overflows. The loose 16-record, 1,024-byte-per-record control is 16,848 wire bytes and is not a lower bound on protection cost.

Cohen's absolute standardized mean difference, recomputed across the 48 constructed plaintext lengths using pooled sample variance, is 12.264 (T1), 5.723 (T2), 3.418 (T3) and 0.033 (T4). Previous values embedded in regime labels were not produced by that calculation and are withdrawn. These descriptive values are not entropy, independent tests or a model of required framing cost.

\paragraph{Adaptive fitting and bootstrap hierarchy}
The clean threshold is fitted to training TLS payload bytes. The defended thresholds use the first four or all twelve training topic pairs. The 1024B defended threshold is 1,680 TLS payload bytes. Stochastic padding has 50 draws; the first is measured on sockets and 49 are calculated using socket-checked record accounting. The independent resampling unit is one of 12 test topics. For each topic, correctness or paired gain is averaged across its two outputs and all draws before 10,000 percentile-bootstrap resamples (seed 20261021). We do not refit the model in that bootstrap or infer uncertainty across fresh model generations. Stored per-topic/per-draw correctness permits exact reconstruction. Both constant caps give a constant specified length sequence without overflow; this says nothing about unobserved timing or metadata.

\paragraph{Retained synthetic tool control}
The original 24 saved inventory/calendar programs are a separate synthetic demonstration. Discard, shadow and read-only routes preserve the prescribed states by construction; 9/24 live-write runs meet both the narrow goal and protected-bystander condition. This is not evidence for autonomous-agent planning or safe deployment of destructive cover.

\subsection{Retail source, permissions and full statistical results}\label{app:tau-cover-details}
The pinned $\tau$-bench commit is \nolinkurl{59a200c6d575d595120f1cb70fea53cef0632f6b}. The instrumented replay preserves the 115 tasks, nine policies, three seeds and two positions (6,210 rows in \path{retail_units.json}); 97 tasks are eligible. The target is the presence of an attempted reference cancellation. Counts include rejected dispatches. The allowed observer receives ordered tool names and response byte lengths, not argument IDs, response contents or backend state. The tool view is separate from the TLS size observer.

The fixed split takes the first 48 sorted eligible task IDs for fitting (7 positive, 41 negative) and the remaining 49 for testing (11 positive, 38 negative). IDs are recorded in \path{retail_diagnostics.json}. Every test task contributes all three seeds. For each task we average accuracy across seeds, then resample 49 tasks 10,000 times for accuracy intervals. Balanced-accuracy intervals independently resample within the 11 positive and 38 negative task strata. These intervals condition on this split, fitted models and saved tasks; policy development was not held out. AUROC uses the uncalibrated ridge linear decision score, including ties. Confusion matrices in \path{retail_observers.json} count 147 task-seed units; they are not 147 independent tasks. Paired $D_2-D_0$ differences are averaged within task before the task bootstrap.

\begin{table*}[t]\centering\footnotesize
\caption{Legacy learned observer on the held-out retail tasks. Intervals are conditional task-bootstrap intervals, not policy-generalization guarantees. The majority baseline is 77.55\% accuracy and 50\% BA.}\label{tab:retail-statistics}
\begin{tabularx}{\textwidth}{@{}Yllll@{}}\toprule
Policy / position & $D_0$ accuracy [95\% CI] & $D_2$ accuracy [95\% CI] & $D_2$ BA [95\% CI] & AUROC\\\midrule
Instruction aware before & 51.02 [36.73, 65.31] & 73.47 [61.22, 85.71] & 47.37 [43.42, 50.00] & 0.634\\
Instruction aware after & 51.02 [36.73, 65.31] & 71.43 [59.18, 83.67] & 52.51 [41.37, 66.15] & 0.694\\
Other-account cancel before & 22.45 [12.24, 34.69] & 71.43 [59.18, 83.67] & 46.05 [40.79, 50.00] & 0.497\\
Other-account cancel after & 22.45 [12.24, 34.69] & 71.43 [59.18, 83.67] & 46.05 [40.79, 50.00] & 0.497\\
Blind cancel before & 22.45 [12.24, 34.69] & 71.43 [59.18, 83.67] & 46.05 [40.79, 50.00] & 0.597\\
Blind cancel after & 22.45 [12.24, 34.69] & 68.03 [54.42, 80.27] & 47.09 [38.16, 58.37] & 0.622\\
\bottomrule\end{tabularx}\end{table*}

The known-position filter scores 100\% accuracy/BA for each policy and position. The cancellation-count rule scores 100\% for blind and other-account cancellation in both positions. Their percentile bootstrap intervals are [100,100] because every observed task is correct; this degeneracy is not a claim about population-level certainty. No identifiers are needed. We do not evaluate an ID-filter attack because IDs are not in the specified view. The legacy model on the other-account after condition recomputes to 71.43\%, correcting the previous prose value 68.03\%; 68.03\% belongs to blind cancellation after insertion.

\paragraph{State metrics and uncertainty}
\path{retail_utility.json} contains task-bootstrap intervals for database equality, full-account equality and equality of reference-changed objects, along with full-census means. The field name \texttt{native\_reward} is retained only for compatibility with the old script; its calculation is orders/users equality, not an invocation of the full reward/response evaluator. Products are shared read-only by the retail tools. The revised account metric checks all orders owned by the active user, plus the user's record. The reference-changed-object set is constructed from baseline differences, not from all accessed objects. Simulator no-cover database equality is a reference-relative 100\%, not proof of user-task success.

\paragraph{Error and denominator reconciliation}
Item self-write before insertion succeeds and changes state in 173/291 units; all remaining 118 errors are unavailable-variant checks. After insertion, 120 errors are non-pending status and 77 are unavailable variants, leaving 94 successful mutations. The raw error strings are retained. Price-difference checks over the 173 initially successful selections find zero differences, but the API still appends a refund and changes the status. The coincidence 204/291 = 68/97 = 70.10\% combines two distinct quantities: preserved reference-changed objects under item self-write and tasks offering an instruction-unmentioned candidate.

For blind cancellation before insertion, database-equal counts by task are: 16:3, 30:3, 31:2, 32:3, 38:3, 54:3, 55:3, 59:2, 66:3, 69:3, 81:2, 99:3. They sum to 33 across 12 tasks. The error-free sensitivity subset is 82 tasks $\times$ 3 seeds = 246 units: 24/246 before and 94/246 after. Simplifying these fractions to 12/123 or 47/123 does not change the underlying sample size. All-unit counts and task IDs are recorded, so repeated seeds are never represented as new independent tasks.

